```latex
\begin{theorem}
The number $\sqrt{7}$ is irrational.
\end{theorem}

\begin{proof}
Suppose, for contradiction, that $\sqrt{7}$ is rational. Then there exist integers $p$ and $q$ such that $q \neq 0$, $\gcd(p,q)=1$, and
\[
\sqrt{7}=\frac{p}{q}.
\]
Squaring both sides gives
\[
7=\frac{p^2}{q^2},
\]
so
\[
p^2=7q^2.
\]
Hence $7 \mid p^2$.

We claim that if $7 \mid n^2$ for an integer $n$, then $7 \mid n$. Indeed, every integer $n$ is congruent modulo $7$ to one of
\[
0,1,2,3,4,5,6.
\]
Squaring these residues modulo $7$ gives
\[
0^2 \equiv 0,\quad
1^2 \equiv 1,\quad
2^2 \equiv 4,\quad
3^2 \equiv 2,\quad
4^2 \equiv 2,\quad
5^2 \equiv 4,\quad
6^2 \equiv 1 \pmod{7}.
\]
Thus the only way for $n^2 \equiv 0 \pmod{7}$ is for $n \equiv 0 \pmod{7}$. Therefore, from $7 \mid p^2$, we get $7 \mid p$.

So there exists an integer $k$ such that
\[
p=7k.
\]
Substituting this into $p^2=7q^2$ gives
\[
(7k)^2=7q^2,
\]
so
\[
49k^2=7q^2.
\]
Dividing by $7$, we obtain
\[
7k^2=q^2.
\]
Therefore $7 \mid q^2$, and by the same argument as above, $7 \mid q$.

Thus $7 \mid p$ and $7 \mid q$, so $p$ and $q$ have a common divisor greater than $1$. This contradicts the assumption that $\gcd(p,q)=1$.

Therefore our original assumption was false, and $\sqrt{7}$ is irrational.
\end{proof}
```